% Editable TikZ: use the force centre, not the centre of an ellipse, for distances.
\providecommand{\ThermalFigure}[3]{\begin{qualityfigure}
\begin{tikzpicture}[>=Latex,font=\small,x=.8cm,y=.8cm]#1\end{tikzpicture}
\figcap{#2}\label{#3}\end{qualityfigure}}
\newcommand{\ThRadialFallBase}[1]{
 \draw[->](-.5,0)--(7,0)node[right]{$r$};\fill(0,0)circle(4pt)node[below]{$M$};
 \fill[DeepBlue](2.5,0)circle(2pt)node[above]{$r_1$};\fill[DeepBlue](6,0)circle(2pt)node[above]{$r_0$};
 \draw[->,thick](5.5,.7)--(3.9,.7)node[midway,above]{下落};#1}
\newcommand{\ThQRadialFall}{\ThermalFigure{\ThRadialFallBase{}}{从外侧静止释放，径向坐标随时间减小，始末距离都为正。}{ch11:quality:radial-q}}
\newcommand{\ThARadialFall}{\ThermalFigure{\ThRadialFallBase{
 \draw[<->](3,-.8)--node[below]{$|\dd r|$}(3.7,-.8);\node[below]at(4.5,-1.6){$\dot r<0,\quad\dd t=-\dd r/|\dot r|$};
}}{径向速度取负分支，时间元为正；积分上下限须与此相配。}{ch11:quality:radial-a}
图~\ref{ch11:quality:radial-a}中位移箭头朝向力心，故能量方程求出的速率须配 $\dot r=-\sqrt{2GM(1/r-1/r_0)}$。时间积分若交换为从 $r_1$ 到 $r_0$，负号才随之消失；这一符号核对避免得到负的下落时间。
}
\newcommand{\ThOrbitBase}[1]{
 \fill[DeepBlue!12](0,0)circle(1.4);\draw[thick](0,0)circle(1.4);\node at(0,0){地球};
 \draw[dashed](0,0)circle(1.65);\fill[DeepBlue](1.65,0)circle(2pt)node[right]{$P$};
 \draw[->](0,0)--(140:1.4)node[midway,above]{$R$};#1}
\newcommand{\ThQOrbit}{\ThermalFigure{\ThOrbitBase{
 \node[right]at(3,1){近地轨道：高度相对 $R$ 可忽略};
}}{圆轨道速度和逃逸速度对应不同的运动要求；图中间隙为示意。}{ch11:quality:orbit-q}}
\newcommand{\ThAOrbit}{\ThermalFigure{\ThOrbitBase{
 \draw[->,thick](1.65,0)--(1.65,1.5)node[right]{$v_c$};
 \draw[->,DeepBlue,thick](1.65,0)--(.5,0)node[above]{$F_g$};
 \node[right,align=left]at(3,0){圆轨道：引力提供向心力\\逃逸临界：无穷远速度趋零};
}}{用 $GM=gR^2$ 把引力参数换为题给地表量；两速度不能混用。}{ch11:quality:orbit-a}
图~\ref{ch11:quality:orbit-a}中的圆周条件是 $mv_c^2/R=GMm/R^2$。逃逸临界态则以无穷远为势能零点并令总机械能为零；二者分别由力与能量决定，故逃逸速度比圆轨道速度大 $\sqrt2$ 倍。
}
\newcommand{\ThApsidesBase}[1]{
 \draw[thick](0,0)ellipse(3 and 1.8);\fill(2.4,0)circle(3pt)node[above left]{$S$};
 \fill(3,0)circle(2pt)node[right]{$P$};\fill(-3,0)circle(2pt)node[left]{$A$};
 \draw[<->](2.4,-2.2)--node[below]{$r_p$}(3,-2.2);\draw[<->](-3,-3)--node[below]{$r_a$}(2.4,-3);#1}
\newcommand{\ThQApsides}{\ThermalFigure{\ThApsidesBase{}}{近点 $P$、远点 $A$ 的距离从焦点力心 $S$ 量起；椭圆中心不作为距离原点。}{ch11:quality:apsides-q}}
\newcommand{\ThTransferBase}[1]{
 \draw[TextGray](0,0)circle(1.2);\draw[TextGray](0,0)circle(3.2);
 \draw[DeepBlue,thick](-1,0)ellipse(2.2 and 1.96);\fill(0,0)circle(3pt)node[below right]{$S$};
 \fill(1.2,0)circle(2pt)node[below right]{$P$};\fill(-3.2,0)circle(2pt)node[left]{$A$};
 \draw[->](0,0)--(50:1.2)node[right]{$r_1$};\draw[->](0,0)--(120:3.2)node[above]{$r_2$};#1}
\newcommand{\ThQTransfer}{\ThermalFigure{\ThTransferBase{}}{转移椭圆与内、外圆轨道分别相切于 $P,A$；力心为它们共同的焦点 $S$。}{ch11:quality:transfer-q}}
\newcommand{\ThATransfer}{\ThermalFigure{\ThTransferBase{
 \draw[->,thick](1.2,0)--(1.2,1.3)node[right]{$\Delta v_1$};
 \draw[->,thick](-3.2,0)--(-3.2,-1)node[left]{$\Delta v_2$};
 \draw[->,DeepBlue,thick](.35,1.55)--(-.55,1.9);
}}{两次脉冲均沿当地顺行切向；无推力段从近点飞到远点，历时半个椭圆周期。}{ch11:quality:transfer-a}
图~\ref{ch11:quality:transfer-a}中的第一脉冲把内圆速度提高到转移椭圆近点速度；第二脉冲把较低的远点速度提高到外圆速度。两处箭头在纸面方向相反，却都沿顺行方向，因此两次速度增量的大小均为正。
}
\newcommand{\ThMassLossBase}[1]{
 \draw[dashed](-1.2,0)circle(1.2);\draw[dashed](-1.2,0)circle(3);
 \fill[DeepBlue](-2.4,0)circle(5pt)node[below]{$M$};\fill(1.8,0)circle(3pt)node[below]{$m$};
 \fill(-1.2,0)circle(1.5pt)node[below]{$C$};\draw[<->](-2.4,-1)--node[below]{$d$}(1.8,-1);
 \draw[->,thick](-2.4,0)--(-2.4,-1.5);\draw[->,thick](1.8,0)--(1.8,1.5);#1}
\newcommand{\ThQMassLoss}{\ThermalFigure{\ThMassLossBase{
 \foreach\a in {30,90,150,210,270,330}{\draw[->,TextGray](-2.4,0)--++(\a:.8);}
 \node[right]at(3,0){$M$ 瞬时各向同性抛出 $\Delta M$};
}}{爆发前绕共同质心 $C$ 转动；无反冲使残余两星瞬时速度连续。}{ch11:quality:mass-loss-q}}
\newcommand{\ThAMassLoss}{\ThermalFigure{
 \fill[DeepBlue](0,0)circle(4pt)node[below]{$M-\Delta M$};\fill(5,0)circle(3pt)node[below]{$m$};
 \draw[<->](0,-.8)--node[below]{$d$（瞬时不变）}(5,-.8);
 \draw[->,thick](5,0)--(5,1.8)node[right]{$v_{\rm rel}$（瞬时不变）};
 \node[above]at(2.5,2.4){相对运动的引力参数改变为 $G(M+m-\Delta M)$};
}{爆发后转入残余双星的相对坐标，束缚条件由相对比能量为负判断。}{ch11:quality:mass-loss-a}
图~\ref{ch11:quality:mass-loss-a}固定的是爆发瞬间的相对距离与相对速度；爆发前的质心位置和分量动能不应继续当作残余系统的质心量。用 $v_{\rm rel}^2=G(M+m)/d$ 与新的引力参数比较，才能直接判断新轨道是否束缚。
}
\newcommand{\ThLRLBase}[1]{
 \fill(0,0)circle(3pt)node[below]{$O$};\draw[->](0,0)--(5,0)node[right]{极轴};
 \fill[DeepBlue](2.4,1.7)circle(2pt)node[right]{$P$};\draw[->](0,0)--(2.4,1.7)node[midway,above left]{$\vect r$};
 \draw(.8,0)arc(0:35.3:.8)node[right]{$\theta$};#1}
\newcommand{\ThQLRL}{\ThermalFigure{\ThLRLBase{
 \draw[->,thick](2.4,1.7)--(1.3,2.5)node[above]{$\vect p$};
}}{平面有心力运动中的位置、动量与极角；待证明守恒矢量可用于选取极轴。}{ch11:quality:lrl-q}}
\newcommand{\ThALRL}{\ThermalFigure{\ThLRLBase{
 \draw[->,DeepBlue,thick](0,0)--(2,0)node[below]{$\vect e$};
 \draw[->,thick](0,0)--(4.8,3.4)node[right]{$\uvec r$ 的方向};
}}{将固定的偏心率矢量作为极轴，点乘径向单位矢量引入 $e\cos\theta$。}{ch11:quality:lrl-a}
图~\ref{ch11:quality:lrl-a}中的 $\theta$ 从常矢量 $\vect e$ 量起，故 $\vect e\cdot\uvec r=e\cos\theta$。另一方面，$\vect L$ 垂直轨道平面；利用三重积把 $(\vect p\times\vect L)\cdot\uvec r$ 化为 $L^2/r$，便把矢量守恒式转成轨道方程。圆轨道 $e=0$ 时极轴可任取。
}
\newcommand{\ThQDust}{\ThermalFigure{
 \draw[thick](0,0)circle(.7);\node at(0,0){$M_s$};\draw[->](0,0)--(-45:.7)node[below]{$r_s$};
 \draw[->,DeepBlue,thick](-4,1.8)--(-1.8,1.8)node[midway,above]{$v$};\draw[dashed](-4,0)--(0,0);
 \draw[<->](-3.7,0)--node[left]{$b$}(-3.7,1.8);\node[right]at(2,1.6){远处平行来流，密度 $\rho$};
}{飞船静止系中的入射尘埃；冲量参数 $b$ 从远处渐近线量起，不等于最近距离。}{ch11:quality:dust-q}}
\newcommand{\ThQPrecession}{\ThermalFigure{
 \draw[thick](0,0)circle(2);\fill(0,0)circle(3pt)node[below]{$M$};\fill[DeepBlue](2,0)circle(2pt)node[above]{$m$};
 \draw[->](2,0)--(.8,0)node[above]{$F_g$};\draw[->,DeepBlue](2,0)--(3,0);
 \node[DeepBlue]at(3.6,.7){$F_1=kmR$};\draw[thin,TextGray](3,0)--(3.4,.45);
 \draw[<->](0,-.7)--node[below]{$R$}(2,-.7);\node[align=left,anchor=west]at(5,1.3){近圆轨道上\\存在微小径向扰动};
}{径向附加力朝外；圆周角运动与径向小振动须分开求频率。}{ch11:quality:precession-q}}
\newcommand{\ThAPrecession}{\ThermalFigure{
 \draw[TextGray,domain=0:360,samples=100]plot({2*cos(\x)/(1+.18*cos(\x))},{2*sin(\x)/(1+.18*cos(\x))});
 \draw[DeepBlue,thick,domain=0:360,samples=100]plot({2*cos(\x+25)/(1+.18*cos(\x))},{2*sin(\x+25)/(1+.18*cos(\x))});
 \fill(0,0)circle(2pt)node[below]{$O$};\draw[dashed](0,0)--(1.695,0);\draw[dashed](0,0)--(25:1.695);
 \draw[->](1,0)arc(0:25:1);\node[right]at(2.8,.5){$\Delta\varpi=2\pi(\Omega/\kappa_r-1)$};
}{两次近日点方向的示意叠图，不把缓慢进动轨道误认为一条闭合椭圆。}{ch11:quality:precession-a}
图~\ref{ch11:quality:precession-a}中一次近日点到下一次近日点用的是径向周期 $2\pi/\kappa_r$，这期间方位角增加 $2\pi\Omega/\kappa_r$。超出 $2\pi$ 的部分就是顺行进动角；因此进动角速度取两频率之差，而不是把附加力简单代入开普勒周期。
}
\newcommand{\ThQOrbitDrag}{\ThermalFigure{
 \draw[thick](0,0)ellipse(3 and 1.8);\fill(2.4,0)circle(3pt)node[below]{$M$};\fill[DeepBlue](0,1.8)circle(2pt)node[above]{$m$};
 \draw[->,thick](0,1.8)--(-1.4,1.8)node[above]{$\vect v$};\draw[->,DeepBlue,thick](0,1.8)--(1.4,1.8)node[above]{$-\gamma\vect v$};
}{弱阻力始终与瞬时速度反向；研究对象是很多周期后的缓慢平均演化。}{ch11:quality:drag-q}}
\newcommand{\ThAOrbitDrag}{\ThermalFigure{
 \draw[thick](-1.8,0)ellipse(3 and 2.4);\draw[DeepBlue,thick](-1.08,0)ellipse(1.8 and 1.44);
 \fill(0,0)circle(3pt)node[below]{$S$};\draw[->](-4,1.6)--(-2.7,1.05);
 \node[right,align=left]at(2,1){同一焦点 $S$\\周期平均下 $a$ 减小、$e$ 保持不变};
}{以力心为中心按比例缩放的相继密切椭圆；实际轨道不在一个周期内突然跳变。}{ch11:quality:drag-a}
图~\ref{ch11:quality:drag-a}中的两条椭圆共用焦点 $S$，两个半轴按同一因子缩小。证明这一点需要同时追踪 $E$ 与 $L$：由它们的平均变化检验 $e^2=1+2EL^2/(G^2M^2m^3)$ 不变，不能只由能量降低就断言偏心率不变。
}
